\documentclass{meetingminutes}
\meetingcommittee{คณะกรรมการประจำคณะเทคโนโลยีอุตสาหกรรม มหาวิทยาลัยราชภัฏอุตรดิตถ์}
\meetingnumber{(เวียนมติ)}{๒๕๖๗}
\meetingdate{วันที่ ๑๑ ธันวาคม พ.ศ. ๒๕๖๗}
\meetingplace{ห้องประชุมคณะเทคโนโลยีอุตสาหกรรม}
\meetingstart{๑๓.๐๐ น.}
\meetingend{๑๖.๐๐ น.}

\begin{document}
\maketitle

\psection{ผู้มาประชุม}
\begin{participants}
  \pmember{1}{ประธานคณะกรรมการประจำคณะเทคโนโลยีอุตสาหกรรม}{คณะเทคโนโลยีอุตสาหกรรม}{ประธาน}
  \pmember{2}{กรรมการประจำคณะ (ผู้แทนคณาจารย์)}{คณะเทคโนโลยีอุตสาหกรรม}{กรรมการ}
  \pmember{3}{ผู้แทนหลักสูตร วศ.ม. วิศวกรรมคอมพิวเตอร์และปัญญาประดิษฐ์}{คณะเทคโนโลยีอุตสาหกรรม}{กรรมการ}
\end{participants}

\psection{ผู้ไม่มาประชุม}
\noindent ไม่มี\par

\startmeeting
\openingchair{ประธานคณะกรรมการประจำคณะ}{กล่าวเปิดการประชุมและดำเนินการประชุมตามระเบียบวาระ ดังนี้}

% ============================================================
\agenda{๑}{เรื่องแจ้งเพื่อทราบ}
% ============================================================

ประธานแจ้งให้ที่ประชุมทราบว่า หลักสูตรวิศวกรรมศาสตรมหาบัณฑิต สาขาวิชาวิศวกรรมคอมพิวเตอร์และปัญญาประดิษฐ์
ได้ผ่านการพิจารณาจากคณะกรรมการบัณฑิตศึกษาประจำมหาวิทยาลัย (๔ มีนาคม ๒๕๖๗) และสภาวิชาการ (๒๑ มีนาคม ๒๕๖๗)
เรียบร้อยแล้ว บัดนี้ให้คณะกรรมการประจำคณะพิจารณาให้ความเห็นชอบเพื่อเสนอต่อสภามหาวิทยาลัยต่อไป

\noindent\textbf{มติที่ประชุม}\quad ที่ประชุมรับทราบ

% ============================================================
\agenda{๒}{พิจารณาให้ความเห็นชอบหลักสูตรวิศวกรรมศาสตรมหาบัณฑิต
สาขาวิชาวิศวกรรมคอมพิวเตอร์และปัญญาประดิษฐ์}
% ============================================================

ผู้แทนหลักสูตรนำเสนอสาระสำคัญของหลักสูตรวิศวกรรมศาสตรมหาบัณฑิต
สาขาวิชาวิศวกรรมคอมพิวเตอร์และปัญญาประดิษฐ์ต่อที่ประชุมคณะกรรมการประจำคณะ
โดยสรุปประเด็นสำคัญดังนี้

\begin{thailist}
  \item ความสอดคล้องกับวิสัยทัศน์ พันธกิจ และยุทธศาสตร์ของคณะเทคโนโลยีอุตสาหกรรม
  \item ผลลัพธ์การเรียนรู้ที่คาดหวัง (PLO) จำนวน ๕ ข้อ ครอบคลุมด้านวิศวกรรมคอมพิวเตอร์และปัญญาประดิษฐ์
  \item โครงสร้างหลักสูตร ๒ แผน ได้แก่ แผน ก (วิทยานิพนธ์) และแผน ข (สารนิพนธ์)
  \item อาจารย์ประจำหลักสูตร ๗ คน มีคุณสมบัติตรงตามเกณฑ์ กมอ. พ.ศ. ๒๕๖๕
\end{thailist}

ที่ประชุมร่วมกันพิจารณา ไม่มีข้อสังเกตเพิ่มเติม

\resolution{ที่ประชุมมีมติให้ความเห็นชอบหลักสูตรวิศวกรรมศาสตรมหาบัณฑิต
สาขาวิชาวิศวกรรมคอมพิวเตอร์และปัญญาประดิษฐ์ คณะเทคโนโลยีอุตสาหกรรม
มหาวิทยาลัยราชภัฏอุตรดิตถ์ และให้เสนอต่อสภามหาวิทยาลัยเพื่ออนุมัติต่อไป}

% ============================================================
\agenda{๓}{เรื่องอื่น ๆ}
% ============================================================

ไม่มีวาระเพิ่มเติม

\noindent\textbf{มติที่ประชุม}\quad ที่ประชุมรับทราบ

\adjournmeeting

\begin{sigblock}
\typist{ผู้แทนหลักสูตร วศ.ม. วิศวกรรมคอมพิวเตอร์และปัญญาประดิษฐ์}
\reviewer{กรรมการประจำคณะ (ผู้แทนคณาจารย์)}{กรรมการ}
\chairsig{ประธานคณะกรรมการประจำคณะ}{ประธานที่ประชุม}
\end{sigblock}

\end{document}
